\section{Introduction}\label{sec: intro}

\subsection{The problem} Given a $3$-manifold~$M$
and a triple~$\Omega = \big(\Omega^1,\Omega^2,\Omega^3\big)$
of closed $2$-forms on~$M$,
it is desired to find a coframing~$\omega = \big(\omega^1,\omega^2,\omega^3\big)$
(i.e., a triple of linearly independent $1$-forms)
satisfying the first-order differential equations
\begin{gather}\label{eq: first-order DEs}
\d\omega^i = \Omega^i
\end{gather}
and the second-order equations that ensure that the metric
\begin{gather}\label{eq: definition of g}
g = \big(\omega^1\big)^2+\big(\omega^2\big)^2+\big(\omega^3\big)^2
\end{gather}
be flat.

This question was originally posed in the context of a problem regarding `residual stress' in elastic bodies due to defects, where the existence of solutions to equations \eqref{eq: first-order DEs} and \eqref{eq: definition of g} is related to the existence of residually stressed bodies that also satisfy a global energy minimization condition. (See~\cite{Acharya} for more details.) However, we feel that the problem is of independent geometric interest.

\subsection{Initial discussion}
As posed, this problem becomes an overdetermined system
of equations for the coframing~$\omega$,
which, in local coordinates~$\big(u^1,u^2,u^3\big)$, can be specified
by choosing the $9$ coefficient functions~$a^i_j(u)$
in the expansion~$\omega^i = a^i_j(u) \d u^j$.
Indeed,~\eqref{eq: first-order DEs}
is a system of $9$ first-order equations
while the flatness of the metric $g$
as defined in~\eqref{eq: definition of g}
is the system of $6$ second-order equations $\Ric(g) = 0$.
Together, these constitute a system of $15$
partial differential equations on the coefficients $a^i_j$
that are independent in the sense that no one of them is a combination of derivatives of the others.

However, the problem can be recast into a different form that makes it more tractable. For simplicity, we will assume that $M$ is connected and simply-connected. The condition that the $\mathbb{R}^3$-valued $1$-form~$\omega$ define a flat metric $g = {}^t\omega \circ \omega$ is then well-known to be equivalent to the condition that~$\omega$
be representable as
\begin{gather*}
\omega = \mba^{-1} \d \mbx,
\end{gather*}
where $\mbx\colon M\to\mathbb{R}^3$ is an immersion and $\mba\colon M\to\SO(3)$ is a smooth mapping.\footnote{In this note, we regard $\mathbb{R}^3$ as \emph{columns}
of real numbers of height $3$, though we will, from time to time, without comment, write them as row vectors in the text.}
This representation is unique up to a replacement of the form
\begin{gather*}
(\mbx,\mba)\mapsto (\mbx',\mba') = (R\mbx + T, R\mba),
\end{gather*}
where $T\in\mathbb{R}^3$ is a constant and $R\in\SO(3)$ is a constant.

Since $\SO(3)$ has dimension~$3$, specifying a pair $(\mbx,\mba)\colon M\to\mathbb{R}^3\times\SO(3)$ is, locally, a choice of $6$ arbitrary (smooth) functions on~$M$.
The remaining conditions on~$\omega$ needed to solve our problem,
\begin{gather}\label{eq: 2_form_eqs}
\d\big(\mba^{-1} \d \mbx\big) = -\mba^{-1} \d \mba \wedge \mba^{-1} \d \mbx = \Omega,
\end{gather}
still constitute $9$ independent first-order equations
for the `unknowns'~$(\mbx,\mba)$ (which are essentially~$6$ in number),
but these equations are not fully independent:
 $\d\Omega=0$ by hypothesis,
and the exterior derivatives of the three $2$-forms
on the left hand side of~\eqref{eq: 2_form_eqs}
also vanish identically for any pair~$(\mbx,\mba)$,
which provides $3$ `compatibility conditions' for the $9$ equations,
thereby, at least formally,
restoring the `balance' of $6$ equations for $6$ unknowns.
Thus, this rough count gives some indication that the problem might
be locally solvable.

However, caution is warranted.
Let $(\bar \mbx,\bar \mba)\colon M\to\mathbb{R}^3\times\SO(3)$
be a~smooth mapping and let~$\bar\Omega = \d\big(\bar \mba^{-1} \d \bar \mbx\big)$.
Linearizing the equations~\eqref{eq: 2_form_eqs}
at the `solution' $(\mbx,\mba) = (\bar \mbx,\bar \mba)$
yields a~system of differential equations of the form
\begin{gather}\label{eq: linearizedeqs}
\d\big( \bar \mba^{-1} (\d \mby - \mbb \d\bar \mbx )\big) = \Psi,
\end{gather}
where \looseness=-1 $(\mby,\mbb)\colon M\to \mathbb{R}^3\oplus \euso(3)$ are unknowns
and $\Psi$ is a closed $2$-form with values in $\mathbb{R}^3$.
If one were expecting~\eqref{eq: 2_form_eqs} to always be solvable,
one might na{\"\i}vely expect \eqref{eq: linearizedeqs} to always
be solvable as well, but this is not so:
When one linearizes at $(\bar \mbx,\bar \mba) = (\bar \mbx, I_3)$,
the linearized system reduces to
\begin{gather}\label{eq: linearizedeqsa=I}
 - \d \mbb\wedge\d\bar \mbx = \Psi,
\end{gather}
where $\mbb\colon M\to\euso(3)\simeq\mathbb{R}^3$
is essentially a set of $3$ unknowns
and $\Psi$ is a given closed $2$-form with values in $\mathbb{R}^3$.
However, as is easily seen, the solvability of~\eqref{eq: linearizedeqsa=I}
for $\mbb$ imposes a system of~$9$ independent first-order
linear equations on~$\Psi$, while the closure of $\Psi$
is only a subsystem of $3$ independent first-order linear equations on~$\Psi$.

Thus, some care needs to be taken in analyzing the system.
Indeed, as Example~\ref{ex: nonsolvability} in~Section~\ref{sec: Zrank1} shows,
there exists an~$\Omega$ defined on a neighborhood of the origin in~$\mathbb{R}^3$
for which there is no solution~$\omega = \mba^{-1} \d \mbx$
to the system~\eqref{eq: 2_form_eqs} on an open neighborhood of the origin.

\subsection{An exterior differential system}
The above observation suggests formulating the problem
as an exterior differential system~$\cI$
on~$X = M\times\mathbb{R}^3\times\SO(3)$
that is generated by the three $2$-form components
of the closed $2$-form
\begin{gather}\label{eq: Theta_defined}
\Theta = -\euroa^{-1} \d \euroa \wedge \euroa^{-1} \d \eurox - \Omega,
\end{gather}
where now, one regards $\eurox\colon X\to\mathbb{R}^3$ and $\euroa\colon X\to\SO(3)$
as projections on the second and third factors.\footnote{We use a different font in equation \eqref{eq: Theta_defined} to emphasize that $\euroa$, $\eurox$, etc., denote matrix- and vector-valued coordinate functions on $X$, while $\mba$, $\mbx$, etc., denote matrix- and vector-valued functions on $M$. We use $\Omega$ to denote both the 2-form on $\mathbb{R}^3$ and its pullback to $X$ via the projection map $\eurox\colon X\to\mathbb{R}^3$.}

We will show that, when~$\Omega$ is suitably nondegenerate,
this exterior differential system is involutive, i.e.,
it possesses Cartan-regular integral flags at every point.
In particular, if~$\Omega$ is also real-analytic,
the Cartan--K\"ahler theorem will imply that the original problem
is locally solvable.

\subsection{Background}
For the basic concepts and results
from the theory of exterior differential systems
that will be needed in this article,
the reader may consult Chapter III of~\cite{BCGGG}. The book \cite{CFB} may also be of interest.

\section{Analysis of the exterior differential system}\label{EDS-sec}

\subsection{Notation}
Define an isomorphism~$[{\cdot}]\colon \mathbb{R}^3\to\euso(3)$
(the space of $3$-by-$3$ skew-symmetric matrices) by the formula
\begin{gather*}%\label{eq: bracketdefined}
\left[ \mbx \right]
= \left[\begin{pmatrix}x^1\\x^2\\x^3\end{pmatrix}\right]
= \begin{pmatrix}
 0 & \phm x^3 & -x^2\\
 -x^3 & 0 & \phm x^1\\
 \phm x^2 & -x^1 & 0
 \end{pmatrix}.
\end{gather*}
The identity $[\mba \mbx] = \mba[\mbx]\mba^{-1}$,
which holds for all $\mba\in\SO(3)$ and $\mbx\in\mathbb{R}^3$,
will be useful, as will the following identities
for $\mbx, \mby\in\mathbb{R}^3$;
$A$ a $3$-by-$3$ matrix with real entries;
$\alpha$ and $\beta$ $1$-forms with values in $\mathbb{R}^3$;
and $\gamma$ a $1$-form with values in $3$-by-$3$ matrices:
\begin{gather}
 [\mbx] \mby = - [\mby] \mbx,\nonumber\\
[A\mbx] = (\tr A) [\mbx] - {}^t\!A [\mbx] - [\mbx] A,\nonumber\\
[\mbx][\mby] = \mby {}^t\mbx - {}^t\mbx \mby I_3,\nonumber\\
[\alpha ]\wedge\beta = [\beta]\wedge\alpha,\nonumber\\
[\gamma\wedge\alpha] = (\tr\gamma)\wedge[\alpha]-{}^t\gamma\wedge[\alpha]+[\alpha]\wedge\gamma,\nonumber\\
[\alpha]\wedge [\beta] = {}^t\beta\wedge\alpha I_3 - \beta\wedge {}^t\alpha,\nonumber\\
{}^t\alpha\wedge[\alpha]\wedge\alpha = - 6 \alpha^1\wedge\alpha^2\wedge\alpha^3,\nonumber\\
[A \alpha]\wedge\alpha = \tfrac12\bigl((\tr A) I_3 - {}^t\!A\bigr) [\alpha]\wedge\alpha.\label{eq: crossproductidentities}
\end{gather}
There is one more identity along these lines that will be useful. It is valid for all $\mathbb{R}^3$-valued $1$-forms~$\alpha$ and functions~$A$ with values in $\GL(3,\mathbb{R})$:
\begin{gather*}
[A\alpha]\wedge A\alpha = \det(A) \big({}^t\!A\big)^{-1} [\alpha]\wedge\alpha.
\end{gather*}

On~$\bbR^3\times\SO(3)$ with first and second factor projections~$\eurox\colon \bbR^3\times\SO(3)\to\bbR^3$ and $\euroa\colon \bbR^3\times\SO(3)\to\SO(3)$,
define the $\mathbb{R}^3$-valued $1$-forms $\xi$ and $\alpha$ by
\begin{gather} \label{basis-forms-on-R3xSO3}
\xi = \euroa^{-1} \d \eurox
\qquad\text{and}\qquad
[\alpha] = \euroa^{-1} \d \euroa = \begin{pmatrix}
0 & \phm\alpha^3 & -\alpha^2\\
-\alpha^3 & 0 & \phm\alpha^1\\
\phm\alpha^2 & -\alpha^1 & 0
\end{pmatrix}.
\end{gather}
These $1$-forms satisfy the so-called `structure equations', i.e., the identities
\begin{gather} \label{structure-eqs-on-R3xSO3}
\d\xi = -[\alpha]\wedge\xi
\qquad\text{and}\qquad
\d\alpha = -\tfrac12 [\alpha]\wedge\alpha.
\end{gather}

\subsection{Formulation as an exterior differential systems problem}
Now suppose that, on~$M^3$, there is specified an $\mathbb{R}^3$-valued, closed $2$-form~$\Omega=\big(\Omega^i\big)$. Choose an $\mathbb{R}^3$-valued coframing~$\eta = (\eta^i)\colon TM\to\bbR^3$.
Then one can write
\begin{gather*}
\Omega = \tfrac12 Z [\eta]\wedge\eta ,
\end{gather*}
where $Z$ is a function on~$M$ with values in $3$-by-$3$ matrices.

Let $\cI$ be the exterior differential system on $X^9 = M\times\mathbb{R}^3\times\SO(3)$ that is generated by the three components of the closed $2$-form
\begin{gather*}%\label{eq: Thetadefined}
\Theta = \d\xi - \Omega = -[\alpha]\wedge\xi - \tfrac12 Z [\eta]\wedge\eta.
\end{gather*}

\begin{Proposition}\label{prop: formulation_as_EDS} If $N^3\subset X$ is an integral manifold of~$\cI$ to which~$\eta$ and~$\xi$ pull back to be coframings, then each point of~$N^3$ has an open neighborhood that can be written as a graph
\begin{gather}\label{eq: xa-graphinX}
\bigl\{\bigl(p,\mbx(p),\mba(p)\bigr)\, \vrule \, p\in U \bigr\}\subset X
\end{gather}
for some open set~$U\subset M$ and smooth maps~$\mbx\colon U\to\mathbb{R}^3$
and $\mba\colon U\to\SO(3)$. Moreover, on~$U$,
the coframing~$\omega = \mba^{-1} \d \mbx$ satisfies~$\d\omega=\Omega$
and the metric~$g = {}^t\omega\circ \omega = {}^t\d \mbx\circ \d \mbx$ is flat.

Conversely, if $U\subset M$ is a simply-connected open subset on which there exists a coframing~$\omega\colon TU\to\bbR^3$ satisfying $(i)$ $\d\omega=\Omega$, and $(ii)$ the metric~$g={}^t\omega\circ \omega$ be flat, then there exist mappings~$\mbx\colon U\to\bbR^3$ and $\mba\colon U\to\SO(3)$ such that $\omega = \mba^{-1} \d \mbx$. Moreover, the immersion $\iota\colon U\to X$ defined by $\iota(p) = \bigl(p,\mbx(p),\mba(p)\bigr)$ is an integral manifold of~$\cI$ that pulls~$\eta$ and~$\xi$ back to be coframings of~$U$.
\end{Proposition}

\begin{proof} The statements in the first paragraph of the proposition are proved by simply unwinding the definitions and can be left to the reader.

For the converse statements (i.e., the second paragraph), suppose that a coframing $\omega\colon TU\to\bbR^3$ be given satisfying the two conditions. By the fundamental lemma of Riemannian geometry, there exists a unique $\bbR^3$-valued $1$-form~$\phi\colon TU\to\bbR^3$ such that
\begin{gather*}
\d\omega = -[\phi]\wedge\omega.
\end{gather*}
The condition that the metric~$g = {}^t\omega\circ \omega$ be flat is then the condition that~$\d\phi = -\tfrac12 [\phi]\wedge\phi$. These equations for the exterior derivatives of $\omega$ and $\phi$, together with the simple-connectivity of~$U$, imply that there exist maps~$\mbx\colon U\to\bbR^3$ and $\mba\colon U\to\SO(3)$ such that
\begin{gather}\label{flat-omega-and-phi}
\omega = \mba^{-1} \d \mbx
\qquad\text{and}\qquad
[\phi] = \mba^{-1} \d \mba.
\end{gather}
Consequently, $g = {}^t\omega\circ \omega$ is equal to ${}^t\d \mbx\circ \d \mbx$, which is flat, by definition. Finally, since~$\d\omega = \Omega$, it follows that the graph manifold~$N^3\subset X$ defined by~\eqref{eq: xa-graphinX} is an integral manifold of~$\cI$. Moreover, since, by construction,
\begin{gather*}
(\mathrm{id}_U, \mbx ,\mba)^*(\xi) = \omega,
\end{gather*}
it follows that $\xi$ and $\eta$ pull back to~$N^3$ to be coframings on~$N^3$.
\end{proof}

\begin{Remark}Observe that the $1$-forms $\omega$ and $\phi$ in equation \eqref{flat-omega-and-phi} are the pullbacks to $U$ of the $1$-forms $\xi$ and $\alpha$, respectively, on $\bbR^3 \times \SO(3)$ defined by equation \eqref{basis-forms-on-R3xSO3}. We will continue to use this notation to distinguish between forms on $\bbR^3 \times \SO(3)$ and their pullbacks via 3-dimensional immersions throughout the paper.
\end{Remark}



\subsection{Integral elements}
By Proposition~\ref{prop: formulation_as_EDS},
proving existence of local solutions of our problem
is equivalent to proving the existence of integral manifolds of~$\cI$
to which $\xi$ and $\eta$ pull back to be coframings. (This latter
condition is usually referred to as an `independence condition'.)

The first step in this approach
is to understand the nature of the integral elements of~$\cI$,
i.e., the candidates for tangent spaces to the integral manifolds of~$\cI$.

A (necessarily $3$-dimensional) integral element~$E\in\mathrm{Gr}(3,TX)$
of~$\cI$ will be said to be \emph{admissible}
if both $\xi\colon E\to\bbR^3$ and $\eta\colon E\to\bbR^3$ are isomorphisms.

\begin{Proposition}\label{prop: K-ordinaryintegralelements}
All of the admissible integral elements of~$\cI$ are K\"ahler-ordinary.\footnote{For definitions of K\"ahler-ordinary, Cartan-ordinary, etc.,
see~\cite[Chapter~III, Definition~1.7]{BCGGG}.}
The set $\cV_3\bigl(\cI,(\xi,\eta)\bigr)$
consisting of admissible integral elements of~$\cI$
is a submanifold of~$\mathrm{Gr}(3,TX)$,
and the basepoint projection~$\cV_3\bigl(\cI,(\xi,\eta)\bigr)\to X$
is a surjective submersion
with all fibers diffeomorphic to~$\GL(3,\bbR)$.
\end{Proposition}

\begin{proof}
Let $(p,\eurox,\euroa)\in X=M\times\bbR^3\times\SO(3)$, and
let $E\subset T_{(p,\eurox,\euroa)} X$ be a $3$-dimensional integral element of $\cI$
to which both $\xi$ and $\eta$ pull back to give an isomorphism of~$E$
with $\mathbb{R}^3$. Then there will exist a $P\in\GL(3,\bbR)$
and a $3$-by-$3$ matrix~$Q$ with real entries
such that $E\subset T_{(p,\eurox,\euroa)} X$ is defined as the kernel of the surjective linear mapping
\begin{gather}\label{eq: Eperp_eqns}
(\xi{-}P \eta, \alpha{-}QP \eta)\colon \ T_{(p,\eurox,\euroa)}\to\bbR^3\oplus\bbR^3.
\end{gather}
To simplify the notation, set $\bar\eta = E^*\eta$.
Then, $E^*\xi = P\bar\eta$ and $E^*\alpha = QP\bar\eta$.
The $2$-form~$\Theta$, which vanishes
when pulled back to~$E$, becomes
\begin{gather*}
0 = E^*\Theta= -[QP \bar\eta]\wedge P \bar\eta-\tfrac12 Z(p) [\bar\eta]\wedge\bar\eta\\
\hphantom{0 = E^*\Theta}{} = -\tfrac12\bigl(\bigl((\tr Q)I_3-{}^t\!Q\bigr)\det(P)\big({}^t\!P\big)^{-1}
 + Z(p)\bigr) [\bar\eta]\wedge\bar\eta.
\end{gather*}
Since~$\bar\eta:E\to\bbR^3$ is an isomorphism, it follows that
\begin{gather*}
\bigl((\tr Q)I_3 - {}^t\!Q\bigr) + Z(p) \,{}^t\!P/\det(P) = 0,
\end{gather*}
so that, solving for~$Q$, one has
\begin{gather}\label{eq: QintermsofPandZ}
Q = \det(P)^{-1}\big(P \, {}^t\!Z(p) - \tfrac12 \tr\bigl(P\, {}^t\!Z(p)\bigr) I_3\big).
\end{gather}

Conversely, if $(p,\eurox,\euroa)\in X=M\times\bbR^3\times\SO(3)$
and~$P\in\GL(3,\bbR)$ are arbitrary
and one defines $Q$ via~\eqref{eq: QintermsofPandZ},
then the kernel~$E\subset T_{(p,\eurox,\euroa)}X$
of the mapping~\eqref{eq: Eperp_eqns}
is an admissible integral element of~$\cI$.

The claims of the Proposition follow directly from these observations.
\end{proof}

\subsection{Polar spaces and Cartan-regularity}
In order to be able to apply the Cartan--K\"ahler theorem to prove existence of solutions in the real-analytic category, one needs a stronger result than Proposition~\ref{prop: K-ordinaryintegralelements}; one needs to show that there are \emph{Cartan}-ordinary admissible integral elements, in other words, to establish the existence of ordinary \emph{flags} terminating in elements of~$\cV_3\bigl(\cI,(\xi,\eta)\bigr)$. This requires some further investigations of the structure of the ideal~$\cI$ near a given integral element in~$\cV_3\bigl(\cI,(\xi,\eta)\bigr)$.

Let $E\in\cV_3\bigl(\cI,(\xi,\eta)\bigr)$ be fixed, with $E\subset T_{(p,\eurox,\euroa)}X$, and let $E$ be defined in this tangent space by the $6$ linear equations
\begin{gather}\label{eq: E in PQterms}
\xi - P \eta = \alpha - QP \eta = 0,
\end{gather}
where $Q$ is given in terms of $P\in\mathrm{GL}(3,\mathbb{R})$ and $Z(p)$ by~\eqref{eq: QintermsofPandZ}. For simplicity, set $\xi_E = (\xi - P \eta)_{\vrule(p,\eurox,\euroa)}$
and $\alpha_E = (\alpha - QP \eta)_{\vrule(p,\eurox,\euroa)}$, and let $\omega_E = (P \eta)_{\vrule(p,\eurox,\euroa)}$. The $9$ components of $\xi_E$, $\alpha_E$, and $\omega_E$ yield a basis of $T^*_{(p,\eurox,\euroa)}X$, with $E^\perp\subset T^*_{(p,\eurox,\euroa)}X$ being spanned by the components of~$\xi_E$ and $\alpha_E$ while $\omega_E\colon E\to\bbR^3$ is an isomorphism.

After calculation using~\eqref{eq: QintermsofPandZ} and the identities~\eqref{eq: crossproductidentities}, one then finds that $\Theta_{\vrule(p,\eurox,\euroa)}$ has the following expression in terms of $\xi_E$, $\alpha_E$, and $\omega_E$:
\begin{gather*}
\Theta_{\vrule(p,\eurox,\euroa)} = - [\alpha_E ]\wedge \omega_E
- [Q \omega_E ]\wedge \xi_E - [\alpha_E ]\wedge \xi_E\\
\hphantom{\Theta_{\vrule(p,\eurox,\euroa)}}{} =-\bigl( [\alpha_E ]+ [\xi_E ]Q\big)\wedge\omega_E
 - \left[\alpha_E\right]\wedge \xi_E .
\end{gather*}
The second term in this final expression, $- [\alpha_E ]\wedge \xi_E$, lies in $\Lambda^2\big(E^\perp\big)$ and hence plays no role in the calculation of the polar equations of~$E$. Hence, the polar spaces for an integral flag of~$E$ can be calculated using only $-\bigl( [\alpha_E ]+ [\xi_E ]Q\big)\wedge\omega_E$.

If $(\mbe_1,\mbe_2,\mbe_3)$ is a basis of~$E$, let $E_i\subset E$ be the subspace spanned by $\{ \mbe_j\, \vrule \, j\le i \}$ and set $w_i = \omega_E(\mbe_i)\in\bbR^3$.
Then the polar space of~$E_i$ is given by
\begin{gather*}
H(E_i) = \big\{\mbv\in T_{(p,\eurox,\euroa)}X\, \vrule\, \big( [\alpha_E(\mbv) ]+ [\xi_E(\mbv) ]Q\big)w_j =0, \ j\le i\big\}.
\end{gather*}
Consequently, the codimension~$c_i$ of this polar space
satisfies $c_i \le 3i$ for $0\le i\le 3$. Since the codimension
of~$\cV_3\bigl(\cI,(\xi,\eta)\bigr)$ in~$\mathrm{Gr}(3,TX)$ is~$9$,
which is always greater than or equal to $c_0 + c_1 + c_2$,
it follows, by Cartan's test, that the flag
$(E_0\subset E_1\subset E_2\subset E_3)$
will be Cartan-ordinary if and only if $c_0+c_1+c_2=9$,
i.e., $c_i = 3i$ for $i = 0,1,2$.
Moreover, this holds if and only if $c_2 = 6$.

Whether or not there is a $2$-plane $E_2\subset E$ with~$c_2 = 6$
evidently depends on~$Q$ (which is determined by~$E$).

\begin{Example}\label{ex: noninvolutivity}
Suppose that $E$ satisfies~$Q = 0$, which, by~\eqref{eq: QintermsofPandZ},
is the case for all of the admissible integral elements based at $(p,\eurox,\euroa)$
if $Z(p) = 0$.
In this case, it is clear that $\left[\alpha_E\right]+\left[\xi_E\right]Q
= \left[\alpha_E\right]$ takes values in skew-symmetric $3$-by-$3$
matrices and hence that, for every $2$-plane $E_2\subset E$,
one must have $H(E_2) = \ker\alpha_E$, so that $c_2=3$.
Thus, Cartan's inequality is strict,
and the integral element~$E$ is not Cartan-ordinary.

Note, though, that this does not imply that there are no solutions
to the original problem on domains containing~$p$ when $Z(p)=0$;
it's just that Cartan--K\"ahler cannot immediately be applied
in such situations. For example, note that,
when~$\Omega$ vanishes identically (equivalently, $Z$ vanishes identically),
then all of the admissible integral elements of~$\cI$
are contained in the integrable $6$-plane field $\alpha = 0$,
and, indeed, the general solution~$\omega$ is of the form~$\omega=\d \mbx$
where $\mbx\colon M\to\bbR^3$ is any immersion.
\end{Example}

For any $3$-by-$3$ matrix $Q$, define $A_Q\subset \eugl(3,\bbR)
= \Hom \big(\bbR^3,\bbR^3\big)$, the \emph{tableau} of~$Q$,
to be the span of the $3$-by-$3$ matrices
\begin{gather*}
[\mbx ]+ [\mby] Q
\end{gather*}
for $\mbx,\mby\in\mathbb{R}^3$.
The dimension of the vector space $A_Q$ lies between $3$ and $6$.

It is evident that the polar equations of flags
in a given admissible integral element $E$
defined by~\eqref{eq: E in PQterms} are governed
by the properties of the tableau~$A_Q$.

To simplify the study of~$A_Q$, it is useful to note that
it has a built-in equivariance:
For $R\in\SO(3)$, one has
\begin{gather*}
R\bigl([\mbx ]+ [\mby] Q\bigr)R^{-1} = R[\mbx]R^{-1}+R[\mby]R^{-1}\big(RQR^{-1}\big) = [R\mbx] + [R\mby] RQR^{-1}.
\end{gather*}
Hence,
\begin{gather*}
R A_Q R^{-1} = A_{RQR^{-1}}.
\end{gather*}
In particular, properties of~$A_Q$ such as its dimension, character sequence, and involutivity depend only on the equivalence class of the matrix~$Q$ under the action of conjugation by $\SO(3)$. Also, writing $Q = q I_3 + Q_0$ where $\tr(Q_0) = 0$, one has
\begin{gather*}
[ \mbx ]+ [\mby ] Q = [ \mbx + q \mby ]+ [\mby ] Q_0 .
\end{gather*}
Thus,
\begin{gather*}
A_Q = A_{Q_0}.
\end{gather*}

\begin{Proposition}\label{prop: nondegenerateA_Q}
The tableau~$A_Q\subset \eugl(3,\bbR) = \Hom \big(\bbR^3,\bbR^3\big)$
has dimension~$6$ and is involutive
with characters~$(s_1,s_2,s_3) = (3,3,0)$,
except when the trace-free part of~$Q$ is conjugate by $\SO(3)$
to a matrix of the form
\begin{gather}\label{eq: Q0normalizeddegenerate}
Q_0 \simeq \begin{pmatrix}-2x&0&0\\0&x+3r&3y\\0&-3y&x-3r\end{pmatrix},
\end{gather}
where $(x,y,r)$ are real numbers satisfying either $r^2 = x^2 + y^2$
or $r = y = 0$.
\end{Proposition}

\begin{proof}
The proof is basically a computation. The conjugation action of $\SO(3)$
on $3$-by-$3$ mat\-rices preserves the splitting of $\eugl(3,\bbR)$
into three pieces: The multiples of the identity (of dimension~$1$),
the subalgebra~$\euso(3)$ (of dimension $3$),
and the traceless symmetric matrices (of dimension $5$).
Moreover, as is well-known, a symmetric $3$-by-$3$ matrix can be
diagonalized by conjugating with an orthogonal matrix. Thus, one
is reduced to studying the case in which $Q_0$ is written in the form
\begin{gather}\label{eq: Q0normalized}
Q_0 = \begin{pmatrix}\phm q_1&\phm p_3&-p_2\\
 -p_3&\phm q_2&\phm p_1\\
 \phm p_2&-p_1&\phm q_3\end{pmatrix},
\end{gather}
where $q_1+q_2+q_3 = 0$.

It is now a straightforward (if somewhat tedious) matter
(which can be eased by MAPLE) to check that, when $A_{Q_0}$
has dimension less than $6$ (the maximum possible),
two of the $p_i$ must vanish. Thus, after conjugating by
a signed permutation matrix that lies in $\SO(3)$,
one can assume that $p_2=p_3=0$.
With this simplification, $A_{Q_0}$ is seen to have dimension
less than $6$ if and only if
\begin{gather*}
p_1 \big({p_1}^2 + 2{q_2}^2 + 5q_2q_3 + 2{q_3}^2\big) =
(q_2 - q_3) \big({p_1}^2 + 2{q_2}^2 + 5q_2q_3 + 2{q_3}^2\big) = 0.
\end{gather*}
Thus, either ${p_1}^2 + 2{q_2}^2 + 5q_2q_3 + 2{q_3}^2=0$ or $p_1=q_2-q_3=0$.
Making the necessary changes of basis,
these two cases give the two non-involutive normal forms
in~\eqref{eq: Q0normalizeddegenerate}.

It remains to show that, when $A_Q$ has dimension~$6$,
it actually is involutive with the stated characters~$(s_1,s_2,s_3)=(3,3,0)$.
To do this, return to the general normal form~\eqref{eq: Q0normalized},
and assume that $A_Q$ has dimension~$6$.
Because $A_Q$ has codimension~$3$ in~$\eugl(3,\bbR)$,
it will be involutive with characters~$(s_1,s_2,s_3)=(3,3,0)$
if and only if it has a non-characteristic covector.
Now, the condition that a covector~$\mbz^* = (z_1,z_2,z_3)\in \big(\bbR^3\big)^*$
be characteristic for~$A_Q$ is the condition
that the $3$-dimensional vector space of rank~$1$ matrices
of the form~$\mbx \mbz^*$ (where $\mbx\in\bbR^3$ and $\mbz^* = (z_1,z_2,z_3)$
is regarded as a row vector) have a nontrivial intersection
with~$A_Q$ in~$\eugl(3,\mathbb{R})$.
The condition that a rank $1$ matrix~$\mathbf{r} = \mathbf{x}\mathbf{z}^*$
lie in the $6$-dimensional subspace $A_Q$ of the $9$-dimensional space
$\frak{gl}(3,\mathbb{R})$ can be expressed as 3 homogeneous linear equations in~$\mathbf{r}$,
i.e., $3$ homogeneous equations bilinear in the components
of~$\mathbf{x}$ and $\mathbf{z}^*$. Regarding $\mathbf{z}^*\not=0$ as given,
this becomes a system of three linear equations for the components
of~$\mathbf{x}$ whose coefficient matrix~$C_Q(\mathbf{z}^*)$ is $3$-by-$3$
with entries that are linear in the components of~$\mathbf{z}^*$.
This system will have a nonzero solution~$\mathbf{x}$
if and only if $\det\bigl(C_Q(\mathbf{z}^*)\bigr) = 0$.
In terms of the coefficients~$p_i$ and~$q_i$ of~$Q_0$,
this determinant vanishing can be written
as a homogeneous cubic polynomial equation
\begin{gather*}%\label{eq: c_Qdefined}
0 = \sum_{ijk} c_{ijk}(p,q) z_i z_j z_k = c_Q(\mbz^*).
\end{gather*}
One then finds (again by a somewhat tedious calculation
that is eased by MAPLE) that this equation holds identically in~$\mbz^*$
(i.e., that all of the~$c_{ijk}(p,q)$ vanish) if and only if~$Q_0$
is equivalent to a matrix of the form~\eqref{eq: Q0normalizeddegenerate}
subject to either of the two conditions~$r=y=0$ or $r^2 = x^2+y^2$.

Thus, except when $Q_0$ is orthogonally equivalent to such matrices,
$A_Q$ has dimension~$6$ and there exists a non-characteristic covector~$\mbz^*$ for~$A_Q$. As already explained, this implies that~$A_Q$ is
involutive, with the claimed Cartan characters.
\end{proof}

\begin{Remark}The $\SO(3)$-orbits of the matrices~$Q$
whose trace-free part~$Q_0$
is of the form~\eqref{eq: Q0normalizeddegenerate}
with~$r=y=0$ forms a closed cone of dimension~$4$
in the ($9$-dimensional) space~$\eugl(3,\mathbb{R})$
of $3$-by-$3$ matrices.
Meanwhile, the $\SO(3)$-orbits of the matrices~$Q$
whose trace-free part~$Q_0$
is of the form~\eqref{eq: Q0normalizeddegenerate}
with~$r^2=x^2+y^2$ forms a closed cone of dimension~$6$
in~$\eugl(3,\mathbb{R})$.

Consequently, the set consisting of those~$Q$
for which $A_Q$ is involutive is an open dense set
in the space~$\eugl(3,\mathbb{R})$.
\end{Remark}

\begin{Remark}\label{rem: Q-hyperbolicity}
It does not appear to be easy to determine the condition on~$Q$
that the real cubic curve~$c_Q(\mbz^*) = 0$ be a smooth,
irreducible cubic with two circuits. This is what one would need
in order to have a chance of showing that the (linearized) equation
were symmetric hyperbolic, which would be a key step in proving
solvability of the original problem in the smooth category.
\end{Remark}

\begin{Corollary}\label{cor: ECartan-ordinary}
If $E\in \cV_3\bigl(\cI,(\xi,\eta)\bigr)$
is defined by equations~\eqref{eq: E in PQterms},
then $E$ is Cartan-regular
if and only if $Q_0 = Q - \tfrac13\tr(Q)I_3$
is not orthogonally equivalent
to a matrix of the form~\eqref{eq: Q0normalizeddegenerate},
where either $r = y =0$ or $r^2 = x^2+y^2$.
\end{Corollary}

\begin{proof}
Everything is clear from Proposition~\ref{prop: nondegenerateA_Q},
except possibly the assertion of Cartan-\emph{regularity}.
However, because the characters are~$(s_1,s_2,s_3)=(3,3,0)$,
when $Q$ avoids the two `degenerate' cones, it follows that,
when $E\in\cV_3\bigl(\cI,(\xi,\eta)\bigr)$ has the property
that its~$A_Q$ is involutive, then, for any non-characteristic
$2$-plane~$E_2\subset E$, we must have $H(E_2) = E$, and hence
$H(E)=E$, so that $E$ must be not only Cartan-ordinary, but
also Cartan-regular.
\end{proof}

\section{Involutivity}\label{sec: Involutivity}
Finally, we collect all of this information together,
yielding our main result:

